% Bibliography copied verbatim from paper/manuscript.tex so that the report and the
% manuscript cite identical records. Regenerate with tools/build_report_data.py if the
% manuscript's bibliography changes.
\chapter*{References}
\addcontentsline{toc}{chapter}{References}
\markboth{References}{}

\begin{thebibliography}{99}
\small

\bibitem{behnke2020} I.~Behnke, L.~Pirl, L.~Thamsen, R.~Danicki, A.~Polze, and O.~Kao,
``Interrupting real-time IoT tasks: How bad can it be to connect your critical embedded system to
the Internet?,'' in \emph{Proc.\ IEEE 39th Int.\ Performance Computing and Communications Conf.\
(IPCCC)}, 2020. doi:10.1109/IPCCC50635.2020.9391536.

\bibitem{behnke2022} I.~Behnke, P.~Wiesner, R.~Danicki, and L.~Thamsen, ``A priority-aware
multiqueue NIC design for real-time IoT devices,'' in \emph{Proc.\ 37th ACM/SIGAPP Symp.\ on
Applied Computing (SAC)}, 2022. doi:10.1145/3477314.3507165.

\bibitem{behnke2023} I.~Behnke, C.~Blumschein, R.~Danicki, P.~Wiesner, L.~Thamsen, and O.~Kao,
``Towards a real-time IoT: Approaches for incoming packet processing in cyber-physical systems,''
\emph{Journal of Systems Architecture}, vol.~140, art.~102891, 2023.
doi:10.1016/j.sysarc.2023.102891.

\bibitem{fischer2020} F.~Fischer \emph{et al.}, ``Analysis of industrial device architectures for
real-time operations under denial of service attacks,'' arXiv:2007.08885, 2020.

\bibitem{timingiso2025} L.~Meszl\'enyi, J.~Kahle, D.~P\"ullen, S.~Kowalewski, S.~Katzenbeisser,
and A.~Kampmann, ``Towards timing isolation for mixed-criticality communication in
software-defined vehicles,'' arXiv:2508.13652, 2025.

\bibitem{mcti2024} D.~Hoornaert, G.~Ghaemi, A.~Bastoni, R.~Mancuso, M.~Caccamo, and G.~Corradi,
``Mcti: Mixed-criticality task-based isolation,'' \emph{Real-Time Systems}, vol.~60, no.~2,
pp.~328--365, 2024. doi:10.1007/s11241-024-09425-5.

\bibitem{mcreview2026} A.~Burns and R.~I.~Davis, ``Mixed criticality systems: A review,'' 14th
edn., Dept.\ of Computer Science, University of York, Feb.\ 2026.

\bibitem{deauther2026} S.~Kremser and K.~Graffi, ``Evaluating the ESP32-S3 for Wi-Fi penetration
testing through the development of Deauther32 and HackHeld32,'' \emph{Sensors}, vol.~26, no.~11,
art.~3287, 2026. doi:10.3390/s26113287.

\bibitem{espwn2023} R.~Cayre, D.~Cauquil, and A.~Francillon, ``ESPwn32: Hacking with ESP32
system-on-chips,'' in \emph{Proc.\ 17th IEEE Workshop on Offensive Technologies (WOOT)}, 2023.

\bibitem{carbajal2026} A.~Carbajal and A.~Jodeiri Akbarfam, ``A software-defined testbed for
quantifying deauthentication resilience in modern Wi-Fi networks,'' arXiv:2602.23513, 2026.

\bibitem{esp32c3re2025} I.~Mudraje, K.~Vogelgesang, J.~Devreker, L.~Gerhorst, P.~Raffeck,
P.~W\"agemann, and T.~Herfet, ``Reverse engineering the ESP32-C3 Wi-Fi drivers for static
worst-case analysis of intermittently-powered systems,'' arXiv:2501.17684, 2025.
doi:10.1145/3722572.3727926.

\bibitem{mqtttls2026} M.~A.~Azhar, M.~Azwar, and O.~Asroni, ``Performance and security analysis
of MQTT over TLS on ESP32-based IoT systems,'' \emph{J.\ Applied Informatics and Computing},
2026.

\bibitem{ondevice2026} V.~Ieropoulos \emph{et al.}, ``Empowering IoT security: On-device
intrusion detection in resource constrained devices,'' arXiv:2605.13159, 2026.

\bibitem{krichen2023} M.~Krichen, ``A survey on formal verification and validation techniques for
Internet of Things,'' \emph{Applied Sciences}, vol.~13, no.~14, art.~8122, 2023.
doi:10.3390/app13148122.

\bibitem{lamport2002} L.~Lamport, \emph{Specifying Systems: The TLA+ Language and Tools for
Hardware and Software Engineers}. Addison-Wesley, 2002.

\bibitem{fuzzsurvey2022} J.~Yun, F.~Rustamov, J.~Kim, and Y.~Shin, ``Fuzzing of embedded
systems: A survey,'' \emph{ACM Computing Surveys}, vol.~55, no.~7, 2022.
doi:10.1145/3538644.

\bibitem{pemu2025} M.~Bley, T.~Scharnowski, S.~W\"orner, M.~Schloegel, and T.~Holz,
``Protocol-aware firmware rehosting for effective fuzzing of embedded network stacks,'' in
\emph{Proc.\ 2025 ACM SIGSAC Conf.\ on Computer and Communications Security (CCS)}, 2025.
doi:10.1145/3719027.3765125.

\bibitem{sysfuss2026} D.~Tharindu, A.~Jayasena, and P.~Mishra, ``SysFuSS: System-level firmware
fuzzing with selective symbolic execution,'' arXiv:2602.02243, 2026.

\bibitem{pirates2024} W.~Alsabbagh, S.~Amogbonjaye, C.~Kim, and P.~Langend\"orfer, ``Pirates of
the MQTT: Raiding IIoT systems with a rogue client,'' in \emph{2024 8th Cyber Security in
Networking Conference (CSNet)}, IEEE, 2024, doi:10.1109/CSNET64211.2024.10851733.

\bibitem{zeghida2024} H.~Zeghida, M.~Boulaiche, and R.~Chikh, ``Security of MQTT protocol: A
brief overview,'' in \emph{Proc. TACC 2024}, CEUR Workshop Proceedings, 2024.

\bibitem{amirkhanova2026} G.~Amirkhanova, S.~Ismailov, A.~Amirkhanov, S.~Adilzhanova,
M.~Zhasuzakova, and S.~Chen, ``A lightweight, end-to-end encrypted data pipeline for IIoT: An
AES-GCM implementation for ESP32, MQTT, and Raspberry Pi,'' \emph{Information}, vol.~17, no.~1,
art.~33, 2026, doi:10.3390/info17010033.

\bibitem{mqttl2026} C.~Yu, M.~Qi, H.~Luo, J.~Tian, and J.~Lu, ``MQTT-L: A lightweight
end-to-end secure protocol with renewable hash chains for IoT environments,'' \emph{Journal of
King Saud University Computer and Information Sciences}, vol.~38, no.~6, art.~539, 2026,
doi:10.1007/s44443-026-00972-4.

\end{thebibliography}
